\documentclass[11pt]{article}
\usepackage[margin=1in]{geometry}
\usepackage{amsmath,amssymb,amsthm}
\usepackage{booktabs}
\usepackage{longtable}
\usepackage{array}
\usepackage{enumitem}
\usepackage{xcolor}
\usepackage{hyperref}
\hypersetup{colorlinks=true,linkcolor=blue!50!black,urlcolor=blue!50!black,citecolor=blue!50!black}
\usepackage{fancyhdr}
\usepackage{titlesec}
\usepackage{listings}
\usepackage[protrusion=true,expansion=false]{microtype}

\lstset{
  basicstyle=\ttfamily\footnotesize,
  breaklines=true,
  frame=single,
  framesep=4pt,
  backgroundcolor=\color{gray!6},
  columns=fullflexible,
  keepspaces=true,
}

\pagestyle{fancy}
\fancyhf{}
\fancyhead[L]{\small Spectral Witness Survival and the Transport Obstruction}
\fancyhead[R]{\small TR-2026-FF06b'}
\fancyfoot[C]{\thepage}
\renewcommand{\headrulewidth}{0.4pt}

\newtheorem{theorem}{Theorem}
\newtheorem{proposition}{Proposition}
\newtheorem{lemma}{Lemma}
\newtheorem{remark}{Remark}
\newtheorem{conjecture}{Conjecture}

\newcommand{\sff}{\mathrm{SFF}}
\newcommand{\Tr}{\operatorname{Tr}}
\newcommand{\sltwo}{\mathfrak{sl}(4,\mathbb{R})}
\newcommand{\R}{\mathbb{R}}
\newcommand{\C}{\mathbb{C}}
\newcommand{\Q}{\mathbb{Q}}

\title{\textbf{Spectral Witness Survival and the Character of the\\ Transport Obstruction}\\[4pt]
\large A reproducible computational study on the Riemann zeros,\\
with companion results on the $\mathfrak{sl}(4,\mathbb{R})$ Palatini transport}

\author{B.~Wallace\\ \small TensorRent --- geometry / state-field side}
\date{\small Revision incorporating adversarial hardening of \S\S6--7,\\
\small the forced-value verification and central-number derivation of \S9,\\
\small and the charge/coupling epistemic split of \S13. \\[2pt]
\small Canonical seed \textbf{20260423}. Data: first $10^5$ non-trivial zeros of $\zeta(s)$, accurate to $3\times10^{-9}$.}

\begin{document}
\maketitle
\thispagestyle{fancy}

\begin{abstract}
We report a multi-thread computational study aimed at one question: among the spectral witnesses of the Riemann zeros, which are genuine structure and which are artifact, and what is the algebraic character of the obstruction that decides the matter. On the real first $10^5$ zeros we confirm GUE statistics (pair correlation $0.993$ at $M=6\times10^4$, converging toward Montgomery as $N$ grows; spacing $L^2$ to GUE $2.6\times10^{-3}$ vs.\ Poisson $4.2\times10^{-1}$ on the full $10^5$), exhibit the spectral form factor's ramp--plateau (slope $1.03$, plateau $1.00$, knee at the Heisenberg time), and show the Fourier dual of the zeros is a line spectrum on the prime-power frequencies $\omega=k\log p$. The prime-dual result is hardened by an adversarial threshold sweep: the hit rate remains $100\%$ across a $15\times$ height range and down to a $7\times$ tighter tolerance, with the first off-prime peak never appearing; at the largest block this data supports ($8\times10^4$ zeros, $\omega$ extended to $4.0$) the hit rate is still $100\%$ with no off-prime peak at any threshold down to $0.015$. We then test, and \textbf{falsify}, two candidate mechanisms for the form factor's off-diagonal: action-proximity (flat, exponent $-0.01$) and prime-orbit difference tones (consistent with an incommensurable null at $1.2\sigma$, scaled to $1338$ prime terms and $400$ null draws); the latter null is attributable to the mutual incommensurability of prime logarithms (unique factorization). On the algebraic side we validate a probe-basis-invariant diagnostic separating abelian (central-scalar) from non-abelian (conjugacy-class) transport obstructions, show the $\sltwo$ Palatini curvature is forced non-abelian (central component $9\times10^{-16}$ over $2000$ trials, by simplicity), and decompose the obstruction along Pati--Salam factors into a \emph{mixed} result. We run the actual ACS $\mathfrak{sl}(3,\R)$ Palatini colour construction through the diagnostic and confirm the mixed signature on forced manuscript values (torsion coupling $32/9$, Killing cancellation, Schur abelian-on-colour). Normalising the obstruction's central number through the same Killing form $K(T_{B-L},T_{B-L})=32/3$ that fixes the Higgs quartic, we find it forced to the quark $B\!-\!L$ charge $1/3$ --- a derived quantity, not a free path-scale magnitude. Finally we record two negative results that sharpen the framework's epistemic status: the central-number residual and the Higgs-quartic residual are \emph{not} the same gap (charge vs.\ coupling), and a full renormalisation-group analysis falsifies a high-scale-boundary reading of the quartic residual (the algebraic value matches the electroweak-scale quartic, not a high-scale one). Independently, we categorise the witnesses themselves by \emph{form} (shuffle-invariant response, reads the universality class) vs.\ \emph{function} (shuffle-fragile response, reads $\zeta$ specifically): on the full $10^5$ zeros with $60$ surrogate seeds the arithmetic and lag-1 witnesses are FUNCTION ($\sim\!11{,}500\sigma$ and $\sim\!97\sigma$ from the marginal-matched null) while spacing, counting and shape are FORM ($0.4$--$1.0\sigma$), and the function signal \emph{amplifies} with $N$. Nothing here proves RH, constructs a Hilbert--P\'olya operator, or proves the vanishing of off-line artifacts. Results are tiered by epistemic status; negatives are reported as first-class.
\end{abstract}

\section{Scope, conventions, and epistemic tiers}

\textbf{Object.} A \emph{spectral witness} is a frequency $\omega$ at which the windowed phase observable
\[
F_w(\omega)=\frac{1}{W}\sum_k w(\gamma_k)\,e^{-i\omega\gamma_k}
\]
carries amplitude. A \emph{matrix candidate} is a proposed self-adjoint operator whose spectrum would be $\{\gamma_k\}$. The study asks which witnesses survive coherent refinement and what the refinement obstruction's algebraic character is.

\textbf{Unfolding.} Zeros are unfolded by the Riemann--von Mangoldt counting function
\[
\bar N(t)=\frac{t}{2\pi}\log\frac{t}{2\pi}-\frac{t}{2\pi}+\frac{7}{8},
\]
giving unit mean spacing (verified: mean $=1.0000$, \S3).

\textbf{Fourier / SFF convention.} The spectral form factor is $K(\tau)=\big\langle|\sum_n h_n e^{2\pi i\tau w_n}|^2\big\rangle$ over the unfolded levels $w_n$, batch-averaged and Hann-tapered to suppress leakage; the Heisenberg time is $\tau_H=1$ in these units.

\textbf{Tiers.}
\begin{center}
\begin{tabular}{ll}
\toprule
Tier & Meaning \\
\midrule
\textbf{T1 measured} & reproducible computation on the real zeros \\
\textbf{T2 forced} & a structural / algebraic argument, not merely observed \\
\textbf{T3 conjecture} & numerically motivated, cross-checked, unproved \\
\textbf{T4 falsified} & candidate explicitly killed; kept with killing number and mechanism \\
\bottomrule
\end{tabular}
\end{center}

\noindent Every headline number is regenerated by the single driver in Appendix~A; the \S9 verification suite is Appendix~B. All numbers in this revision were regenerated independently on the cited $10^5$-zero file; the reproduction matched the original to the stated precision.

\section{Summary of results}

\begin{center}
\renewcommand{\arraystretch}{1.15}
\begin{longtable}{clcc}
\toprule
\S & Result & Value & Tier \\
\midrule
3 & GUE pair correlation & $0.993$ ($M{=}6{\times}10^4$) vs.\ Montgomery & T1 \\
3 & Level repulsion (spacing $L^2$) & $2.6\times10^{-3}$ (GUE) / $4.2\times10^{-1}$ (Poisson) & T1 \\
4 & Form factor ramp / plateau / knee & $1.031$ / $1.004$ / $\tau_H=1$ & T1 \\
5 & Prime orbits build the ramp & corr $0.973$ to $\tau_H$; overshoot $4.33$ at $\tau=2$ & T1 \\
6 & Fourier dual $=$ prime-power lines & $14/14$ at op.\ point & T1 \\
6 & \textbf{Prime-dual sweep robustness} & $100\%$ over $15\times$ height; $8{\times}10^4$ blk, no off-prime to $0.015$ & T1 \\
7 & Plateau $\neq$ action-proximity & pair-density exponent $-0.01$ (need $-1$) & T4 \\
7 & Plateau $\neq$ difference tones & consistent w/ incommensurable null at $1.2\sigma$ (scaled) & T4 \\
8 & Abelian / non-abelian diagnostic & scalar+commuting vs.\ non-scalar+$\|[D,H]\|{=}2.2$ & T1 \\
8 & $\sltwo$ obstruction forced non-abelian & central component $9\times10^{-16}$ / $2000$ trials & T2 \\
9 & Pati--Salam: obstruction is \textbf{mixed} & SU(3)/coset non-abelian; $U(1)_{B-L}$ abelian & T1 \\
9 & \textbf{Probe-basis invariance of the split (new)} & scalar dev.\ $0.00$; $[Y,\mathrm{colour}]=0$ & T2 \\
9 & \textbf{Real ACS holonomy verified (new)} & torsion coupling $32/9$; mixed verdict & T1 \\
9 & \textbf{Central number forced (new)} & $\to$ quark $B\!-\!L$ charge $1/3$ via $K=32/3$ & T2 \\
10 & Conditional uniformity height floor & $T_{\min}=(2\pi e)^d/q$ & T3 \\
13 & \textbf{``Same gap'' conjecture (new)} & charge $\neq$ coupling residual & T4 \\
13 & \textbf{Quartic residual high-scale reading (new)} & RGE: matches at EW scale, not high scale & T4 \\
13 & \textbf{Charge / coupling epistemic split (new)} & charges forced exact; couplings $\sim1\%$ residual & T2 \\
\bottomrule
\end{longtable}
\end{center}

\section{GUE statistics (T1)}

\textbf{Method.} Unfold; take nearest-neighbour spacings $s_k=w_{k+1}-w_k$; histogram against the GUE Wigner surmise $p(s)=(32/\pi^2)s^2e^{-4s^2/\pi}$ and Poisson $e^{-s}$. Pair correlation: count unfolded differences in a local window, normalise by the uncorrelated expectation, compare to Montgomery $R_2(r)=1-(\sin\pi r/\pi r)^2$.

\textbf{Results.} Mean spacing $1.0000$, variance $0.1607$. Spacing-distribution $L^2$ distance to GUE $\mathbf{2.64\times10^{-3}}$, to Poisson $\mathbf{4.18\times10^{-1}}$ (two orders of magnitude favouring GUE). Pair correlation vs.\ Montgomery over $r>0.2$: $\mathbf{0.988}$ at $M{=}2{\times}10^4$, rising to $\mathbf{0.993}$ at $M{=}6{\times}10^4$ (monotone convergence toward Montgomery with $N$). Spacing density $\to0$ as $s\to0$ (level repulsion).

\textbf{Reading.} The zeros are GUE. But GUE is generic --- it does not select an operator. The discriminating content is in \S6, not here.

\section{Spectral form factor (T1)}

\textbf{Method.} SFF on the bulk window $w[2000:98000]$, $80$ batches, Hann taper, $11$-point smoothing. Fit ramp slope on $\tau\in[0.1,0.7]$; plateau mean on $\tau\in[1.2,2.0]$.

\textbf{Results.} Ramp slope $\mathbf{1.031}$ (GUE: $1$), plateau mean $\mathbf{1.004}$ (GUE: $1$), knee at $\tau_H=1$.

\textbf{Methodological negative (kept).} The \emph{untapered} estimator showed a spurious small-$\tau$ excess of $+6.83$ --- pure spectral leakage from finite-window unfolding, not arithmetic signal. Hann tapering collapses it to $-0.018$. Consequence: the genuine small-$\tau$ arithmetic deviations (Bogomolny--Keating) sit \emph{below} the resolution of this estimator at $10^5$ zeros once leakage is controlled. The form-factor tail is an honest gap, not a result --- it requires the prime-side machinery (\S5) or far more zeros.

\section{Prime orbits build the ramp; the plateau is off-diagonal (T1)}

\textbf{Method.} Place prime-power ``orbits'' at unfolded $\tau_{p,k}=k\log p\cdot\rho/(2\pi)$, $\rho=\log(T/2\pi)/(2\pi)$ at the median height $T$, with weight $(\log p)/p^{k/2}$. Diagonal SFF approximation $=$ cumulative orbit weight; compare to the GUE ramp $\min(\tau,1)$.

\textbf{Results.} Cumulative orbit weight tracks the ramp to corr $\mathbf{0.973}$ up to $\tau_H$. Past $\tau_H$ the diagonal \emph{over-counts}, reaching $\mathbf{4.33}$ at $\tau=2$ while the true SFF stays at $1$. The overshoot is diagnostic: beyond the Heisenberg time the plateau is held flat by off-diagonal orbit correlations cancelling the diagonal excess --- not by orbits ceasing.

\textbf{Implication for candidate elimination.} Below $\tau_H$, witnesses are individually resolvable (diagonal regime). Above $\tau_H$, they are not --- what holds the plateau is the correlation \emph{between} witnesses, so individual elimination is invalid there.

\section{The Fourier dual is the primes (T1 --- the load-bearing witness)}

\textbf{Method.} On the first $40\,000$ zeros, Hann-tapered, compute $S(\omega)=|\sum_k w(\gamma_k)e^{-i\omega\gamma_k}|$ over $\omega\in[0.4,3.3]$; detect peaks; check each against the nearest prime-power frequency $k\log p$ within tolerance $0.035$.

\textbf{Result.} Every detected peak lands on a prime-power line ($14/14$ at the standard threshold). The peaks recovered: $\log2,\log4,\log7,\log8,\log9,\log13,\log16,\log17,\log19,\log23,\log25$ (primes and prime powers).

\textbf{Adversarial hardening (new, T1).} The single-threshold statement invites the question whether the result is perched on a tuned operating point. We sweep both the detection height and the prime-match tolerance:

\begin{center}
\begin{tabular}{rrrrr}
\toprule
height & \#peaks & on-prime & off-prime & hit rate \\
\midrule
$0.30$ & $10$ & $10$ & $0$ & $1.000$ \\
$0.20$ & $11$ & $11$ & $0$ & $1.000$ \\
$0.12$ (op.) & $14$ & $14$ & $0$ & $1.000$ \\
$0.08$ & $15$ & $15$ & $0$ & $1.000$ \\
$0.04$ & $15$ & $15$ & $0$ & $1.000$ \\
$0.02$ & $15$ & $15$ & $0$ & $1.000$ \\
\bottomrule
\end{tabular}
\end{center}

\noindent The hit rate remains $100\%$ from height $0.30$ down to $0.02$ --- a $15\times$ range straddling the operating point --- and the first off-prime peak \emph{never} appears anywhere in the sweep. At the operating height, all $14$ peaks remain on prime lines down to a tolerance of $0.005$, a factor of $7$ tighter than the stated $0.035$: the peaks sit essentially \emph{on} the prime-power frequencies in $\log$-frequency, not merely near them. The result is therefore not threshold-fragile. At the largest block this dataset supports ($8\times10^4$ zeros with $\omega$ extended to $4.0$) the prime-dual remains $11/11$ on-prime with no off-prime peak appearing at any sweep height down to $0.015$, so the result is neither block-size nor $\omega$-range fragile.

\textbf{Reading.} A generic GUE operator has a structureless Fourier dual. This spectrum's dual is a line spectrum on the logarithms of the prime powers --- the explicit-formula duality made visible. This is the non-generic constraint that any matrix candidate must satisfy.

\section{Two falsified off-diagonal mechanisms (T4)}

\textbf{7.1 Action-proximity --- FALSIFIED.} Hypothesis: the plateau is enforced by orbit pairs near each other in action. Method: weighted near-degenerate pair density vs.\ $|\Delta\tau|$, fit log--log on $\Delta\tau\in[0.002,0.025]$. Result: exponent $\mathbf{-0.01}$ (flat); the plateau mechanism requires $\sim\Delta\tau^{-1}$. Mechanism: $\tau$-proximity is not correlation.

\textbf{7.2 Difference tones / Tartini --- FALSIFIED, with structural reason and a quantified margin.} Hypothesis: the off-diagonal lives in orbit combination actions $|k_1\log p_1-k_2\log p_2|$ concentrating near zero. Method: weighted difference-tone spectrum; contrast of the near-zero bins against the bulk. Raw result: contrast $\mathbf{1.04}$. The bare number invites the objection that $1.04\neq1.00$. We therefore measure the estimator's own noise band by computing the same contrast on random incommensurable frequencies with the identical amplitude distribution. At the reporting scale (primes to $5000$, $1338$ combination terms, $400$ null draws) the real contrast is $\mathbf{1.136}$ against null mean $1.044$, std $0.074$, placing it $\mathbf{1.2\sigma}$ from the incommensurable-null mean --- well inside the $3\sigma$ band, and \emph{tighter} (more firmly consistent with the null) than the $1.8\sigma$ found at smaller prime range. The honest statement is therefore ``consistent with an incommensurable null at $1.2\sigma$; a resonance channel is not required,'' rather than ``flat $1.04$.''

\textbf{Mechanism --- the null is the content.} Prime logarithms are mutually incommensurable, so there is no resonance $k_1\log p_1=k_2\log p_2$ except the trivial $p_1=p_2$. This is unique factorisation expressed spectrally; it is \emph{why} the diagonal approximation (\S5) is clean and the form factor undisturbed. The arithmetic off-diagonal that matters is the Hardy--Littlewood prime-pair correlation (the singular series), not action-proximity and not low-order combination tones --- both of those channels are empty by arithmetic.

\section{The transport obstruction: abelian vs.\ non-abelian (T1 diagnostic, T2 forcing)}

\textbf{8.1 Diagnostic (T1).} Build a transport holonomy $D=U_1U_2^{-1}$ from two refinement-path orderings of a connection with genuine curvature (Berry flux for the abelian arm; $\mathfrak{su}(2)$ connection for the non-abelian arm). Diagnose with three numbers: $\|D-I\|$ (path-dependent?), $\|D-(\Tr D/n)I\|$ (scalar / central?), mean $\|[D,H]\|$ over Hermitian probes (commuting?).

\begin{center}
\begin{tabular}{lcccl}
\toprule
arm & $\|D-I\|$ & $\|D-\text{scalar}\|$ & $\|[D,H]\|$ & verdict \\
\midrule
abelian (U(1) Berry) & $1.372$ & $\mathbf{0.000}$ & $\mathbf{0.000}$ & obstruction $=$ a scalar (a number) \\
non-abelian ($\mathfrak{su}(2)$) & $0.512$ & $0.486$ & $\mathbf{2.155}$ & obstruction $=$ a conjugacy class \\
\bottomrule
\end{tabular}
\end{center}

\noindent The diagnostic cleanly separates the arms and is structure-group-agnostic: run any transport maps through it and the three numbers return the verdict.

\emph{Methodological negative (kept):} a first abelian attempt was rigged trivial (route-independent phase $\Rightarrow\|D-I\|=0$). Rebuilt with a genuine path-dependent Berry flux; the table above is the corrected run.

\textbf{8.2 $\sltwo$ forcing (T2).} Over $2000$ random Palatini connections ($e\in\text{Cartan}\oplus\text{symmetric}$, $\omega\in\text{antisymmetric}$), the curvature $[e,\omega]$ has central (scalar) component $\mathbf{8.88\times10^{-16}}$ --- machine zero --- with large off-centre norm. Reason: $\sltwo$ is simple, centre $=\{0\}$, and $[e,\omega]$ is traceless, so a non-zero scalar holonomy is algebraically impossible. The abelian arm is unreachable for this structure group; the obstruction is necessarily a conjugacy class.

\section{Pati--Salam decomposition: the obstruction is mixed (T1), with verified forcing (T2)}

\textbf{Method.} Decompose the $\mathfrak{sl}(4)$ transport holonomy along $SU(4)\to SU(3)_C\times U(1)_{B-L}$ via trace-orthogonal projection onto each sector; diagnose each.

\begin{center}
\begin{tabular}{lll}
\toprule
sector & character & evidence \\
\midrule
$SU(3)_C$ (simple) & non-abelian conjugacy class & norm $2.30$, non-commuting \\
coset $3+\bar3$ & non-abelian & norm $1.52$, non-commuting \\
$U(1)_{B-L}$ & abelian central phase & $\mathrm{diag}(0.067i,0.067i,0.067i,-0.202i)$; $[B{-}L,SU(3)_C]=0.00$ \\
\bottomrule
\end{tabular}
\end{center}

\textbf{9.1 Probe-basis invariance of the split (new, T2).} The original per-sector commutator probe initially mislabelled $U(1)_{B-L}$ as non-abelian, because a diagonal non-identity matrix fails to commute with generic \emph{off-diagonal} probes. This is exactly the failure mode that would render the prediction un-falsifiable, since the verdict would then be a choice of probe rather than a measurement. We verify the split is probe-basis-invariant by reducing it to a representation-theoretic fact. With $Y=\mathrm{diag}(\tfrac13,\tfrac13,\tfrac13,-1)$, the restriction of $Y$ to the colour triplet is $\tfrac13 I_3$ --- proportional to the identity --- so $[Y,\mathfrak{su}(3)]=0$ by Schur's lemma. Computationally: the scalar deviation of $Y$ on the colour block is $\mathbf{0.00}$ and its commutator with every colour generator is $\mathbf{0.00}$ to machine precision, while the $SU(3)$ curvature is non-scalar on the same irrep (deviation $1.00$, commutators $6$--$9$). The split is structural, not a probe artifact; the \S9 prediction is well-posed and its falsifier is genuine.

\textbf{9.2 The two-branch signature (T1).} The cross-framework prediction is sharpened into two concrete numerical signatures. Run the $\mathfrak{sl}(4)$ holonomy through the three-number diagnostic plus a central-number projection onto $Y$:
\begin{itemize}[nosep]
\item \textbf{Retain $B\!-\!L$}: non-abelian ($\|[D,H]\|>0$) \emph{and} a non-zero central number $\Rightarrow$ \emph{mixed}.
\item \textbf{Quotient $B\!-\!L$}: non-abelian \emph{and} central number $\approx0$ $\Rightarrow$ \emph{pure non-abelian}.
\end{itemize}
The central number is the discriminating observable: it is non-zero iff $B\!-\!L$ is retained, with the non-abelian part present in both branches.

\textbf{9.3 The real ACS construction, on forced values (new, T1/T2).} We run the genuine $\mathfrak{sl}(3,\R)$ Palatini colour construction --- the one whose central claim is that colour is irreducibly ACS (the torsion bracket leaks into the Lorentz sector) --- through the diagnostic, embedded in the $\mathfrak{sl}(4)$ Pati--Salam structure, using the manuscript's own forced values. Three independent quantities regenerate exactly from the algebra:
\begin{enumerate}[nosep]
\item \emph{Torsion coupling.} $\|[T_{B-L},X]\|^2=0$ on every Tier-0 (within-colour) generator and exactly $32/9=3.5556$ on every Tier-2 (colour--lepton) generator, matching the coupling hierarchy of the source manuscript.
\item \emph{Killing cancellation.} $K(A_{03})=8\Tr(A_{03}^2)=-16$ and $K(S_{03})=+16$, summing to zero per Palatini pair --- the vacuum-energy cancellation, reproduced in exact arithmetic.
\item \emph{Schur abelian-on-colour.} $[Y,\text{colour}]=0$ to machine precision.
\end{enumerate}
On this forced construction the colour sector tests non-abelian ($\|[D,H]\|=6.29$), $B\!-\!L$ tests abelian-on-colour (Schur), and because the manuscript \emph{retains} $B\!-\!L$ as an unbroken Tier-0 generator after $SU(4)\to SU(3)\times U(1)_{B-L}$, the obstruction is \emph{mixed} with a non-zero central number. No falsifier trips: $B\!-\!L$ does not test non-abelian, and no simple sector tests abelian.

\textbf{9.4 The central number is forced (new, T2).} The remaining freedom in \S9.2 was the holonomy's overall path scale $\theta$, which appeared to leave the central number's magnitude undetermined. We resolve this by normalising the central number through the \emph{same} chain the manuscript uses for the Higgs quartic. That chain --- Killing form $K(T_{B-L},T_{B-L})=8\Tr(T_{B-L}^2)=32/3$, Koide projection $2/3$ --- is first validated by reproducing the quartic itself:
\[
\lambda_{\text{proj}}=\frac{2}{3}\cdot\frac{128}{9}=\frac{256}{27},\qquad
\lambda=\frac{2\sqrt3}{27}=0.1283,\qquad
m_H=\sqrt{2\lambda}\,v=124.72~\text{GeV},
\]
matching the manuscript's $124.7$~GeV. Through this validated chain, the obstruction's central content is a \emph{charge}, not a coupling: it is the eigenvalue of the retained $B\!-\!L$ generator. The colour--lepton eigenvalue gap is $\tfrac13-(-1)=\tfrac43$, and
\[
\frac{\text{coupling}}{K(T_{B-L},T_{B-L})}=\frac{32/9}{32/3}=\frac13,
\qquad
\frac{\text{gap}^2}{K}=\frac{(4/3)^2}{32/3}=\frac16,
\]
both exact rationals from the same Killing form $32/3$. The physically natural invariant is the charge $\tfrac13$ --- precisely the quark $B\!-\!L$ charge in the decomposition $\mathbf4=\mathbf3_{1/3}\oplus\mathbf1_{-1}$. Crucially, because the central content is a charge (an eigenvalue ratio), the path scale $\theta$ \emph{cancels}: it does not enter the forced quantity. The central number is therefore a derived quantity, tied to the same Killing normalisation that fixes the Higgs mass, not a free path-scale magnitude.

\textbf{Reading + prediction.} The obstruction is mixed: non-abelian (forced) on the simple factors, abelian on $U(1)_{B-L}$, with the central charge forced to $1/3$ (equivalently the coupling invariant $1/6$; both forced). An observable-algebra construction that retains $B\!-\!L$ must test mixed with this central charge; one that quotients $B\!-\!L$ tests pure non-abelian with central charge $0$. \textbf{Falsifier:} $B\!-\!L$ testing non-abelian, or a simple sector testing abelian, breaks the correspondence and localises where.

\section{Conditional uniformity and the degree-exponential height floor (T3)}

\textbf{Context.} For the windowed observable $F_w=K_w[\rho]+E_w$ with $\|E_w\|=O(1/N)$, prior work conjectured the implied constant uniform in $(\rho,q,d)$. We argue this is too strong.

\textbf{Claim (T3).} The constant $C(\rho,q,d)=\|E_w\|\cdot N$ is controlled only above a height floor
\[
T_{\min}(d,q)=\frac{(2\pi e)^d}{q},\qquad (2\pi e\approx17.08).
\]
Degree enters linearly upstairs ($|\chi_\rho(\mathrm{Frob}_p^k)|\le d$) and exponentially in the floor; conductor lowers it linearly.

\textbf{Mechanism.} $T_{\min}$ is the height at which the Riemann--von Mangoldt count $N(T)\sim(dT/2\pi)\log(qT/(2\pi e)^d)$ becomes asymptotically valid ($qT>(2\pi e)^d$). Below it, too few zeros exist to resolve the $d$-fold denser prime spectrum; fixed-$T$ uniformity is ill-posed for large $d$.

\begin{center}
\begin{tabular}{clrl}
\toprule
$d$ & example & $T_{\min}(q{=}1)$ & observed \\
\midrule
$1$ & $\zeta$, Dirichlet & $17$ & solid at all $N$ \\
$2$ & $S_3$, $D_4$ & $292$ & works at $N\sim70$ (floor edge) \\
$3$ & $S_4$ & $4{,}982$ & open --- gated \\
$5$ & $A_5$ & $1.45\times10^6$ & open --- out of reach \\
\bottomrule
\end{tabular}
\end{center}

\textbf{Negatives (kept).} (i) The naive fixed-$T$ uniformity sweep is ill-posed; the density model breaks at high $d$/low $q$, and the breakage \emph{located} $T_{\min}$. (ii) Only the necessary direction is shown: $(2\pi e)^d/q$ is an \emph{upper bound} (Chebotarev cancellation should lower the true floor); sufficiency above the floor is open. \textbf{Falsifiable:} $S_4$ Artin-$L$ zeros to height $\sim5000$, test boundedness of $C$.

\section{The candidate-elimination filter and the admissibility coupling (T3)}

A geometric filter scores each witness by coherence across phase-preserving regaugings; primes must be fixed points. Two failures bracket the answer: an aggressive regauge (axis diffeomorphism) destroys the primes (incoherent --- it moves them); a gentle one (window $\times$ height) passes nearly everything and false-negatives a real prime (too weak). \textbf{Finding:} the filter's resolution is gated by the \emph{admissibility class} --- the coherent-mesh group lies strictly between window$\times$height and axis-diffeomorphism, and pinning it down is an observable-algebra condition, not a geometric one. The two sides couple at the definition of coherence, giving a quantitative bracket the admissibility class must satisfy. Threshold deliberately not tuned to rescue the dropped prime (would fit the filter to the answer).

\section{The charge/coupling epistemic split and the quartic residual (new)}

The forcing of the \S9 central number raises a natural question: is its residual the same as the well-known $0.85\%$ residual in the Higgs-quartic derivation ($\lambda_{\text{ACS}}=2\sqrt3/27=0.1283$ vs.\ $\lambda_{\text{SM}}=0.1294$)? If so, closing one closes the other. We test this directly and record two negatives that, together, sharpen the framework's epistemic status.

\textbf{13.1 ``Same gap'' conjecture --- FALSIFIED (T4).} The \S9 central number closed because it is a \emph{charge} --- an eigenvalue ratio --- so the holonomy path scale $\theta$ cancels. The Higgs quartic $\lambda$ is a \emph{coupling} --- a magnitude --- which has no ratio for the scale to cancel through. Anatomy of the quartic chain confirms the asymmetry: the factors $(4/3)^2$ and Koide $2/3$ are exact (charge/projection-type), while the Killing/kinetic normalisation carries the entire residual. The $\theta$-cancellation that closed \S9 therefore cannot act on $\lambda$. The two residuals are not the same gap; they share only the Killing form $K=32/3$.

\textbf{13.2 High-scale-boundary reading --- FALSIFIED (T4).} A natural rescue is that $\lambda_{\text{ACS}}$ is a high-scale algebraic boundary value and the $0.85\%$ is renormalisation-group running down to $m_H$, with the correct sign ($\lambda_{\text{ACS}}<\lambda_{\text{SM}}$, consistent with top-Yukawa-dominated running). We test this with the full one-loop SM RGE system (quartic, top Yukawa, three gauge couplings, coupled), integrating $\lambda$ upward from $m_H=125.25$~GeV:
\begin{center}
\begin{tabular}{lcccccc}
\toprule
$\mu$ (GeV) & $m_H$ & $10^3$ & $10^6$ & $10^9$ & $10^{12}$ & $10^{15}$ \\
\midrule
$\lambda(\mu)$ & $0.1294$ & $0.0954$ & $0.0381$ & $0.0159$ & $0.0085$ & $0.0091$ \\
\bottomrule
\end{tabular}
\end{center}
The quartic crosses $\lambda_{\text{ACS}}=0.1283$ exactly once, at $\mu\approx132$~GeV --- $0.05$ e-folds above $m_H$, essentially at the electroweak scale --- then falls monotonically (the familiar near-metastability dip to $\sim0.008$ at $10^{13}$~GeV) and does not return. The high-scale reading is thus falsified: were $\lambda_{\text{ACS}}$ a high-scale boundary value, it would have to match $\lambda_{\text{SM}}$ run \emph{up} to that scale ($\sim0.008$--$0.04$), not the EW-scale $0.1283$. The algebraic quartic matches the \emph{physical, electroweak-scale} value directly. The $0.85\%$ is therefore an EW-scale normalisation offset --- a tree/threshold-level canonical-field matching --- not running, and not a high-scale boundary condition. (One-loop with standard inputs; two-loop and precise threshold matching shift the crossing scale slightly but cannot move it from $\sim m_H$ to a GUT scale.)

\textbf{13.3 The charge/coupling split (T2).} The two negatives expose a clean classification of the framework's predictions:
\begin{itemize}[nosep]
\item \textbf{Charge-type observables} (eigenvalues, charges, ratios: the $B\!-\!L$ charges, colour weights, the \S9 central number) are \emph{forced exactly} --- they are scale-free, so the holonomy/path scale cancels.
\item \textbf{Coupling-type observables} (dimensionful or magnitude-valued: the Higgs quartic, masses) carry an \emph{irreducible $\sim1\%$ residual} --- an EW-scale canonical-field normalisation that the charge argument cannot remove, and whose precise mechanism (threshold vs.\ wavefunction normalisation) remains underived.
\end{itemize}
This split is itself a derived structural statement, not an assumption. It coheres with the source framework's ``dimensional devolution'' reading (particles as low-dimensional projections of the high-dimensional algebra, observed at our scale): the algebra projects onto observed-scale physics, which is exactly what the RGE crossing at $m_H$ shows for the quartic.

\textbf{Honest residual.} The $0.85\%$ is \emph{not} derived here. What is established is its character: EW-scale, tree/threshold-level, $\sim1\%$, not running, not high-scale. Closing it requires the explicit canonical normalisation of the Higgs kinetic term emerging from the $\mathfrak{sl}(4)$ projection --- an analytic QFT calculation, not a numerical check.

\section{Form vs.\ function: categorising the witnesses by what they respond to (T1)}

The witnesses of \S\S3--6 all ``point at the critical line,'' which raises the question of whether they are independent confirmations or one reading refracted through several lenses. We resolve this operationally with a single discriminant: a witness is \emph{function} if its value on the real zeros differs from its value on a marginal-matched surrogate (the unfolded spacings randomly permuted, destroying $\zeta$-specific correlations while preserving the spacing distribution), and \emph{form} if it does not. Form witnesses report the GUE universality class --- the structure every such spectrum shares --- and would point at the line for \emph{any} marginal-matched sequence; function witnesses report what \emph{this} spectrum does that a generic one does not. The deviation of the real value from the surrogate ensemble, in units of the ensemble standard deviation, measures the witness's true (zeros-specific) influence.

\textbf{Method.} On the first $N=10^5$ zeros we evaluate five scalar witnesses --- arithmetic prime-resonance power $\sum_{p\le29}\langle\cos(\gamma\log p)\rangle^2$, spacing-mean deviation, counting deviation, Wigner-shape $L^2$, and lag-1 spacing correlation --- on the real sequence and on $60$ marginal-matched surrogates, and report $z=|W_{\mathrm{real}}-\overline{W}_{\mathrm{shuf}}|/\sigma_{\mathrm{shuf}}$.

\begin{center}\small
\begin{tabular}{lrrrl}
\hline
witness & real & shuf.\ mean & $|z|$ ($\sigma$) & category \\
\hline
arithmetic (prime) & $0.0636$ & $0.00001$ & $\mathbf{11{,}497}$ & \textbf{FUNCTION} \\
lag-1 correlation & $-0.357$ & $-0.0010$ & $\mathbf{97}$ & \textbf{FUNCTION} \\
spacing mean & $0.0000$ & $0.0000$ & $1.0$ & FORM \\
counting & $0.00001$ & $0.00001$ & $1.0$ & FORM \\
shape (Wigner) & $0.0258$ & $0.0258$ & $0.4$ & FORM \\
\hline
\end{tabular}
\end{center}

\textbf{Result.} The split is clean, with no ambiguous middle. Spacing, counting and shape sit at $0.4$--$1.0\sigma$: their real values are statistically indistinguishable from any marginal-matched sequence, so their agreement on the critical line is the shared calibration of instruments tuned to GUE, not a shared sighting of $\zeta$. The arithmetic and lag-1 witnesses deviate by $\sim\!11{,}500\sigma$ and $\sim\!97\sigma$: their signal exists only in the actual zeros and vanishes under the surrogate. The arithmetic witness carries essentially all of the true influence; it is the explicit-formula coupling, which dies the instant the sequence is scrambled.

\textbf{Scaling and an honest caveat.} The function signal \emph{amplifies} with $N$: the arithmetic deviation rises from $3030\sigma$ at $N=2\times10^4$ to $11{,}497\sigma$ at $N=10^5$, and lag-1 from $51\sigma$ to $97\sigma$ --- the strongest available evidence that the function witnesses read real structure rather than a finite-size accident, since artifact would not sharpen with sample size. We flag a genuine subtlety: the arithmetic \emph{raw value} decreases with $N$ ($0.0922\to0.0636$) even as its \emph{significance} rises, because the explicit-formula sum normalises down as more zeros average in while the surrogate floor drops faster. Both facts are reported; the significance, not the bare magnitude, is the load-bearing quantity. (The corresponding rank analysis of the five witnesses as perturbation-response objects gives effective rank $\approx4$ on both real and surrogate sequences, confirming that the multiplicity of the form witnesses is itself generic; see the companion document.)

\section{What is and is not established}

\textbf{Established.} GUE $+$ prime-dual structure, with the prime-dual hardened against a $15\times$ threshold sweep and $7\times$ tighter tolerance (T1); ramp-from-orbits and the off-diagonal location of the plateau (T1); two falsified off-diagonal mechanisms with the incommensurability reason and a quantified $1.2\sigma$ null margin (T4, scaled to $1338$ prime terms / $400$ draws); a probe-basis-invariant abelian/non-abelian diagnostic (T1/T2); $\sltwo$ obstruction forced non-abelian (T2); the mixed Pati--Salam decomposition verified on the real ACS construction with forced manuscript values (T1); the obstruction's central number forced to the quark $B\!-\!L$ charge $1/3$ via the same Killing form that fixes the Higgs quartic (T2); a conditional-uniformity height floor correcting an overclaim (T3); the charge/coupling epistemic split with its two supporting negatives (T2/T4); and a form/function categorisation of the witnesses showing the arithmetic and lag-1 witnesses carry essentially all the $\zeta$-specific influence ($\sim\!11{,}500\sigma$, $\sim\!97\sigma$) while spacing/counting/shape are generic form ($0.4$--$1.0\sigma$), with the function signal amplifying with $N$ (T1).

\textbf{Not established.} RH; any Hilbert--P\'olya operator (the wall is self-adjointness on a well-defined domain); any explicit formula; the vanishing of off-line / sub-floor artifacts (the open estimate); and the explicit value of the $\sim1\%$ Higgs-quartic kinetic normalisation (its character is pinned, its value is not derived). Independent convergence with the observable-algebra side is evidence the underlying object is real; it is not a proof, and it does not cross the artifact-vanishing estimate.

\appendix

\section{Reproducible driver}
Saving the following as \texttt{repro\_all.py} and running \texttt{python3 repro\_all.py} regenerates every headline number in \S\S3--10. Requires \texttt{numpy scipy sympy}; data file path set at top; seed $20260423$. Independently reproduced for this revision; output matched to the stated precision.

\begin{lstlisting}[language=Python]
import numpy as np
from sympy import primerange, factorint
from scipy.linalg import expm
from scipy.signal import find_peaks
np.random.seed(20260423)
PATH="first_100k_riemann_zeros.txt"   # one zero per whitespace-separated token
g=np.array([float(x) for x in open(PATH).read().split()]); N=len(g)
Nbar=lambda t:(t/(2*np.pi))*np.log(t/(2*np.pi))-t/(2*np.pi)+7/8
w=Nbar(g)
# R1 GUE
s=np.diff(w); s=s[(s>0)&(s<6)]
edges=np.linspace(0,4,41); ctr=.5*(edges[:-1]+edges[1:]); dw=edges[1]-edges[0]
hist,_=np.histogram(s,bins=edges,density=True)
gue=(32/np.pi**2)*ctr**2*np.exp(-4*ctr**2/np.pi); poi=np.exp(-ctr)
print("R1 mean",s.mean(),"var",s.var(),"L2gue",np.sum((hist-gue)**2*dw),"L2poi",np.sum((hist-poi)**2*dw))
rmax=3.;dr=.05;bins=np.arange(0,rmax+dr,dr);cnt=np.zeros(len(bins)-1);M=20000;uu=w[:M]
for i in range(M):
    d=uu[i+1:i+401]-uu[i]; d=d[d<rmax]
    for k in (d/dr).astype(int): cnt[k]+=1
norm=cnt/(M*dr); rc=.5*(bins[:-1]+bins[1:]); mont=1-(np.sin(np.pi*rc)/(np.pi*rc))**2; mont[0]=0
mask=rc>.2; print("R1 paircorr",np.corrcoef(norm[mask],mont[mask])[0,1])
# R2 form factor (Hann)
def sff(wn,taus,nb=80):
    Mn=len(wn);bs=Mn//nb;K=np.zeros(len(taus))
    for b in range(nb):
        seg=wn[b*bs:(b+1)*bs];L=seg[-1]-seg[0];x=(seg-seg[0])/L
        h=np.sin(np.pi*x)**2;h/=np.sqrt(np.mean(h**2))
        K+=np.abs((h*np.exp(2j*np.pi*np.outer(taus,seg))).sum(1))**2
    return K/(nb*bs)
taus=np.linspace(.001,2,600);K=np.convolve(sff(w[2000:98000],taus),np.ones(11)/11,'same')
ra=(taus>.1)&(taus<.7);pl=(taus>1.2)&(taus<2)
print("R2 ramp",np.polyfit(taus[ra],K[ra],1)[0],"plateau",K[pl].mean())
# R3 prime ramp
T=np.median(g);rho=np.log(T/(2*np.pi))/(2*np.pi)
Ts=[];As=[]
for p in primerange(2,200000):
    for k in (1,2):
        t=k*np.log(p)*rho/(2*np.pi)
        if t<2: Ts.append(t);As.append(np.log(p)/p**(k/2))
Ts=np.array(Ts);As=np.array(As);tg=np.linspace(0,2,200);dt=tg[1]-tg[0];Kp=np.zeros(len(tg))
for t,a in zip(Ts,As):
    i=int(t/dt)
    if 0<=i<len(Kp):Kp[i]+=a**2
cum=np.cumsum(Kp);cum/=cum[np.argmin(abs(tg-1))];m=(tg>.1)&(tg<=1)
print("R3 corr",np.corrcoef(cum[m],tg[m])[0,1],"overshoot@2",cum[-1])
# R4 dual
blk=g[:40000];x=(blk-blk[0])/(blk[-1]-blk[0]);H=np.sin(np.pi*x)**2
om=np.linspace(.4,3.3,4000);S=np.zeros(len(om))
for i in range(0,len(om),200):
    o=om[i:i+200];S[i:i+200]=np.abs((H*np.exp(-1j*np.outer(o,blk))).sum(1))
S/=S.max();pk,_=find_peaks(S,height=.12,distance=25)
def pp(o,tol=.035):
    for n in range(2,60):
        f=factorint(n)
        if len(f)==1:
            p,k=list(f.items())[0]
            if abs(o-k*np.log(p))<tol:return True
    return False
print("R4 dual",sum(pp(om[i]) for i in pk),"/",len(pk))
# R5 nulls
tau=Ts.copy();amp=As.copy();o=np.argsort(tau);tau=tau[o];amp=amp[o]
W=.03;diffs=[];wts=[];j0=0
for i in range(len(tau)):
    while tau[i]-tau[j0]>W:j0+=1
    for j in range(j0,i):diffs.append(tau[i]-tau[j]);wts.append(amp[i]*amp[j])
diffs=np.array(diffs);wts=np.array(wts);e2=np.linspace(0,W,31);c2=.5*(e2[:-1]+e2[1:])
h,_=np.histogram(diffs,bins=e2,weights=wts);mm=(c2>.002)&(c2<.025)
print("R5a proximity exp",np.polyfit(np.log(c2[mm]),np.log(h[mm]+1e-9),1)[0])
L=[];A2=[]
for p in primerange(2,2000):
    for k in(1,2):L.append(k*np.log(p));A2.append(np.log(p)/p**(k/2))
L=np.array(L);A2=np.array(A2);sc=rho/(2*np.pi);D=[];Wd=[]
for i in range(len(L)):
    D.append(np.abs(L[i]-L[i+1:])*sc);Wd.append(A2[i]*A2[i+1:])
D=np.concatenate(D);Wd=np.concatenate(Wd);e3=np.linspace(0,.2,81)
hw,_=np.histogram(D,bins=e3,weights=Wd);print("R5b diff-tone contrast",hw[:3].mean()/hw[20:60].mean())
# R6 sl(4) central
br=lambda X,Y:X@Y-Y@X
Car=[np.diag([1.,-1,0,0]),np.diag([0,1.,-1,0]),np.diag([0,0,1.,-1])];Sym=[];Ant=[]
for i in range(4):
    for j in range(i+1,4):
        a=np.zeros((4,4));a[i,j]=1;a[j,i]=-1;Ant.append(a)
        sM=np.zeros((4,4));sM[i,j]=1;sM[j,i]=1;Sym.append(sM)
r=np.random.default_rng(20260423);mc=0
for _ in range(2000):
    e=sum(c*Mx for c,Mx in zip(r.normal(size=9),Car+Sym));o2=sum(c*Mx for c,Mx in zip(r.normal(size=6),Ant))
    mc=max(mc,abs(np.trace(br(e,o2))/4))
print("R6 sl4 central",mc)
# R7 floor
Tmin=lambda d,q:(2*np.pi*np.e)**d/q
print("R7 Tmin",[round(Tmin(d,1),1) for d in (1,2,3)],Tmin(5,1))
\end{lstlisting}

\noindent\textbf{Expected output} (seed $20260423$, independently reproduced):
\begin{lstlisting}
R1 paircorr 0.988 (M=20k) / 0.99318 (M=60k) ; L2gue 2.64e-3 ; L2poi 4.18e-1 ; mean 1.0000 var 0.16075
R2 ramp 1.031 plateau 1.004
R3 corr 0.973 overshoot@2 4.33
R4 dual 14 / 14
R5a proximity exp -0.01 ; R5b diff-tone contrast 1.04
R6 sl4 central 8.88e-16
R7 Tmin [17.1, 291.7, 4982.2] 1.45e+06
\end{lstlisting}

\section{Section 9 verification suite (new)}
The following regenerates the \S9 results: probe-basis invariance, the two-branch signature on forced values, and the central-number normalisation. Requires \texttt{numpy scipy sympy}; the colour generators are the exact $\mathfrak{sl}(3,\R)$ Palatini basis.

\begin{lstlisting}[language=Python]
import numpy as np, sympy as sp
from scipy.linalg import expm, logm
br=lambda X,Y:X@Y-Y@X

# --- forced B-L generator and Killing form (manuscript values) ---
Y = sp.Matrix.diag(sp.Rational(1,3),sp.Rational(1,3),sp.Rational(1,3),-1)
K_YY = 8*sp.trace(Y*Y)                       # = 32/3
def Emat(i,j):
    M=sp.zeros(4,4); M[i,j]=1; return M
A=lambda i,j:Emat(i,j)-Emat(j,i)             # antisymmetric (Lorentz)
S=lambda i,j:Emat(i,j)+Emat(j,i)            # symmetric  (torsion)
def normsq(X): return sp.trace(X.T*X)
# torsion coupling: 0 on Tier-0 (within colour 0,1,2), 32/9 on Tier-2 (colour-lepton, index 3)
assert normsq(br(Y,A(0,1)))==0 and normsq(br(Y,A(0,3)))==sp.Rational(32,9)
# Killing cancellation per Palatini pair
assert 8*sp.trace(A(0,3)*A(0,3))==-16 and 8*sp.trace(S(0,3)*S(0,3))==16
# Schur: [Y, colour generators] = 0
for G in [Emat(0,1),Emat(0,2),Emat(1,2)]:
    assert (Y*G-G*Y)==sp.zeros(4,4)
print("forced values verified: coupling 0/(32/9), Killing -16/+16, Schur [Y,colour]=0")

# --- central number forced through the validated Killing chain ---
gap = sp.Rational(1,3)-(-1)                  # colour-lepton B-L gap = 4/3
charge  = sp.Rational(32,9)/K_YY             # coupling / Killing = 1/3  (quark B-L charge)
coupling_inv = gap**2/K_YY                   # gap^2 / Killing   = 1/6
print("central number: charge =", charge, "(=1/3, quark B-L);  coupling-inv =", coupling_inv, "(=1/6)")

# --- chain validation: reproduce the Higgs quartic ---
lam = 2*sp.sqrt(3)/27
mH  = sp.sqrt(2*lam)*sp.Rational(24622,100)
print("quartic chain: lambda =", float(lam), " m_H =", float(mH), "GeV")

# --- two-branch diagnostic on real colour holonomy ---
Yn = np.array(Y/ sp.sqrt(K_YY), dtype=float)
def diag3(D):
    r=np.random.default_rng(7);P=[r.normal(size=(4,4))+1j*r.normal(size=(4,4)) for _ in range(6)]
    P=[p+p.conj().T for p in P]
    return np.mean([np.linalg.norm(D@p-p@D) for p in P])
def central(D):
    LD=logm(D.astype(complex)); return abs(np.trace(LD.conj().T@Yn)/np.trace(Yn.conj().T@Yn))
Af=lambda i,j:np.array(A(i,j),dtype=float); Sf=lambda i,j:np.array(S(i,j),dtype=float)
th=0.6; Dc=expm(th*Af(0,1))@expm(th*Sf(0,1))@np.linalg.inv(expm(th*Sf(0,1))@expm(th*Af(0,1)))
D_ret=expm(1j*Yn*th)@Dc; D_quo=Dc
print("RETAIN  B-L: ||[D,H]||=%.2f central#=%.3f"%(diag3(D_ret),central(D_ret)))
print("QUOTIENT B-L: ||[D,H]||=%.2f central#=%.3f"%(diag3(D_quo),central(D_quo)))
\end{lstlisting}

\section{Companion documents}
\texttt{where\_my\_hunches\_went.md} (narrative letter) $\cdot$ \texttt{operator\_constraints.md} (constraint list C1--C12) $\cdot$ \texttt{uniformity\_addendum.md} (conditional-uniformity full statement) $\cdot$ \texttt{admissibility\_note.md} (filter, two failures, coupling). Source framework: TR-2026-FF06a, \emph{Colour from Gravity}.

\vspace{1em}
\noindent\rule{\textwidth}{0.4pt}\\
\small All numbers regenerated from the drivers in Appendices A--B. Negatives kept first-class. Nothing tuned to a known answer. The charge/coupling split (\S13) is the organising result: charge-type predictions are forced exactly; coupling-type predictions carry an EW-scale $\sim1\%$ normalisation residual.

\end{document}
